\uuid{lQDv}
\titre{Module/argument}
\auteur{bodin ; hanani ; mzari}
\organisation{exo7}
\licence{CC BY-NC-SA 4.0}
\datecreate{2020-01-20}
\niveau{L1}
\module{Algèbre}
\chapitre{Nombres complexes}
\sousChapitre{Forme cartésienne, forme polaire}
\difficulte{1}
\theme{Écritures algébrique et géométrique, module, argument}
\exo7qcmid{lille-1.nombres-complexes.004}

\contenu{
\begin{question}
\qclassification{104.01}
\qtags{motcle=module/argument}
Soit $\theta \in \Rr$. $e^{i\theta}\in \Rr$  si et seulement si :
\begin{answers}
    \bad{ $\theta  =0$ }

    \bad{ $\theta  =2\pi$}

    \bad{$\theta  = 2k\pi$, $k \in \Zz$}

    \good{ $\theta  =k\pi$, $k \in \Zz$}
\end{answers}
\begin{explanations}
$e^{i\theta}= \cos \theta + i \sin \theta $ et $\sin \theta = 0 $ si et seulement si $\theta  =k\pi$, $k \in \Zz$.
\end{explanations}
\end{question}
}
